% --- Archon physics formalization source begin ---
% archon:physics
% archon:covers PhyXMiniProblems/problem_phyx_mini_0697.lean
% archon:source-report reports/phyx_mini/problem_phyx_mini_0697.source.json

\chapter{Physics problem phyx\_mini\_0697}
\label{ch:PhyXMiniProblems_problem_phyx_mini_0697}

\paragraph{Problem source.}
\#\# Physical scenario

Find the center of mass of a thin, uniform rod of length \textbackslash{}(L\textbackslash{}) and mass \textbackslash{}(M\textbackslash{}). Use this result. A \textbackslash{}(1.60\textbackslash{},\textbackslash{}mathrm\{m\}\textbackslash{})-long rod rotates about its center of mass with an angular acceleration of \textbackslash{}(6.0\textbackslash{},\textbackslash{}mathrm\{rad/s\textasciicircum{}2\}\textbackslash{}).

\#\# Question

Find the tangential acceleration of the rod.

\#\# Answer choices

- A: A: \$3.8\textbackslash{}

- B: B: \$5.8\textbackslash{}

- C: C: \$4.8\textbackslash{}

- D: D: \$4.3\textbackslash{}

\#\# Figure caption (auxiliary; use the image as primary evidence)

The image depicts a rod on a Cartesian coordinate plane with a horizontal \textbackslash{}(x\textbackslash{})-axis and a vertical \textbackslash{}(y\textbackslash{})-axis. The rod extends from the origin \textbackslash{}((0,0)\textbackslash{}) to a length \textbackslash{}(L\textbackslash{}) along the \textbackslash{}(x\textbackslash{})-axis, labeled as "Rod."

There is a small segment of the rod highlighted, labeled as a "small cell of width \textbackslash{}(dx\textbackslash{})" at position \textbackslash{}(x\textbackslash{}). The width of this segment is marked as \textbackslash{}(dx\textbackslash{}).

The image includes a dotted arrow pointing to the highlighted segment with the text: 

"A small cell of width \textbackslash{}(dx\textbackslash{}) at position \textbackslash{}(x\textbackslash{}) has mass \textbackslash{}(dm = (M/L)dx\textbackslash{})."

Here, \textbackslash{}(M\textbackslash{}) represents the total mass of the rod, and \textbackslash{}(L\textbackslash{}) is its length. The formula describes the mass \textbackslash{}(dm\textbackslash{}) of the small cell in terms of the rod's density \textbackslash{}(M/L\textbackslash{}) and its width \textbackslash{}(dx\textbackslash{}).

\#\# Dataset metadata

Rotational Motion; Physical Model Grounding Reasoning, Numerical Reasoning

\paragraph{Recorded answer/context.}
C: C: \$4.8\textbackslash{}

\paragraph{Figure/image path.}
/root/proposal\_for\_physic/science-mango/phyx\_mini\_run/phyx\_data/test\_image/697.png

\paragraph{Formalization target.}
create a compiling Lean file with sorry bodies at `PhyXMiniProblems/problem\_phyx\_mini\_0697.lean`. The Lean declarations must preserve the physical quantities, dimensions or dimensional roles, figure labels, governing-law hypotheses, and final relation expressed by this problem.
Use Mathlib/Physlib names found through LeanExplore where available. If a physics API is missing, introduce faithful local abstractions rather than scalar placeholder aliases.

\begin{theorem}[Physics formalization target]
\label{thm:physics:phyx_mini_0697:target}
\lean{PhyXMiniProblems.ProblemPhyXMini0697.problem_phyx_mini_0697}
\uses{def:physics:phyx-mini-0697:phyxminiproblems-problemphyxmini0697-lengthinmeters, def:physics:phyx-mini-0697:phyxminiproblems-problemphyxmini0697-accelerationinmeterspersecondsquared, def:physics:phyx-mini-0697:phyxminiproblems-problemphyxmini0697-rotatingrodexperiment, def:physics:phyx-mini-0697:phyxminiproblems-problemphyxmini0697-matchesprimaryrodfigure, def:physics:phyx-mini-0697:phyxminiproblems-problemphyxmini0697-obeysthinuniformrodmassmodel, def:physics:phyx-mini-0697:phyxminiproblems-problemphyxmini0697-matchesrotatingrodproblem, def:physics:phyx-mini-0697:phyxminiproblems-problemphyxmini0697-obeysrigidrotationkinematics, def:physics:phyx-mini-0697:phyxminiproblems-problemphyxmini0697-haspositionindependenttangentialacceleration, def:physics:phyx-mini-0697:phyxminiproblems-problemphyxmini0697-recordedanswerchoice, def:physics:phyx-mini-0697:phyxminiproblems-problemphyxmini0697-matchesdisplayedtangentialaccelerationatposition}
The assigned autoformalize agent should translate this physics problem into Lean declarations in the covered file, with theorem and lemma proof bodies written as `by sorry`.
\end{theorem}
\begin{proof}
This is an autoformalization task, not a proof task. Produce statements that can later be proved by the physics prover without weakening the physical model.
\end{proof}

% --- Archon declaration topology begin ---
\section{Declaration topology}
\begin{definition}[Mass Quantity]
  \label{def:physics:phyx-mini-0697:phyxminiproblems-problemphyxmini0697-massquantity}
  \lean{PhyXMiniProblems.ProblemPhyXMini0697.MassQuantity}
  A nonnegative physical mass
\end{definition}
\begin{proof}
  This is the declared model definition of Mass Quantity.
\end{proof}
\begin{definition}[Length Quantity]
  \label{def:physics:phyx-mini-0697:phyxminiproblems-problemphyxmini0697-lengthquantity}
  \lean{PhyXMiniProblems.ProblemPhyXMini0697.LengthQuantity}
  A nonnegative physical length
\end{definition}
\begin{proof}
  This is the declared model definition of Length Quantity.
\end{proof}
\begin{definition}[Linear Mass Density Quantity]
  \label{def:physics:phyx-mini-0697:phyxminiproblems-problemphyxmini0697-linearmassdensityquantity}
  \lean{PhyXMiniProblems.ProblemPhyXMini0697.LinearMassDensityQuantity}
  A nonnegative mass per unit length
\end{definition}
\begin{proof}
  This is the declared model definition of Linear Mass Density Quantity.
\end{proof}
\begin{definition}[Angular Acceleration Magnitude Quantity]
  \label{def:physics:phyx-mini-0697:phyxminiproblems-problemphyxmini0697-angularaccelerationmagnitudequantity}
  \lean{PhyXMiniProblems.ProblemPhyXMini0697.AngularAccelerationMagnitudeQuantity}
  An angular-acceleration magnitude; radians are dimensionless
\end{definition}
\begin{proof}
  This is the declared model definition of Angular Acceleration Magnitude Quantity.
\end{proof}
\begin{definition}[Acceleration Magnitude Quantity]
  \label{def:physics:phyx-mini-0697:phyxminiproblems-problemphyxmini0697-accelerationmagnitudequantity}
  \lean{PhyXMiniProblems.ProblemPhyXMini0697.AccelerationMagnitudeQuantity}
  A nonnegative tangential-acceleration magnitude
\end{definition}
\begin{proof}
  This is the declared model definition of Acceleration Magnitude Quantity.
\end{proof}
\begin{definition}[mass In Kilograms]
  \label{def:physics:phyx-mini-0697:phyxminiproblems-problemphyxmini0697-massinkilograms}
  \lean{PhyXMiniProblems.ProblemPhyXMini0697.massInKilograms}
  \uses{def:physics:phyx-mini-0697:phyxminiproblems-problemphyxmini0697-massquantity}
  Kilogram readout of a physical mass
\end{definition}
\begin{proof}
  This is the declared model definition of mass In Kilograms.
\end{proof}
\begin{definition}[length In Meters]
  \label{def:physics:phyx-mini-0697:phyxminiproblems-problemphyxmini0697-lengthinmeters}
  \lean{PhyXMiniProblems.ProblemPhyXMini0697.lengthInMeters}
  \uses{def:physics:phyx-mini-0697:phyxminiproblems-problemphyxmini0697-lengthquantity}
  Metre readout of a physical length
\end{definition}
\begin{proof}
  This is the declared model definition of length In Meters.
\end{proof}
\begin{definition}[linear Mass Density In Kilograms Per Meter]
  \label{def:physics:phyx-mini-0697:phyxminiproblems-problemphyxmini0697-linearmassdensityinkilogramspermeter}
  \lean{PhyXMiniProblems.ProblemPhyXMini0697.linearMassDensityInKilogramsPerMeter}
  \uses{def:physics:phyx-mini-0697:phyxminiproblems-problemphyxmini0697-linearmassdensityquantity}
  Kilogram-per-metre readout of a linear mass density
\end{definition}
\begin{proof}
  This is the declared model definition of linear Mass Density In Kilograms Per Meter.
\end{proof}
\begin{definition}[angular Acceleration In Radians Per Second Squared]
  \label{def:physics:phyx-mini-0697:phyxminiproblems-problemphyxmini0697-angularaccelerationinradianspersecondsquared}
  \lean{PhyXMiniProblems.ProblemPhyXMini0697.angularAccelerationInRadiansPerSecondSquared}
  \uses{def:physics:phyx-mini-0697:phyxminiproblems-problemphyxmini0697-angularaccelerationmagnitudequantity}
  Radian-per-second-squared readout of angular-acceleration magnitude
\end{definition}
\begin{proof}
  This is the declared model definition of angular Acceleration In Radians Per Second Squared.
\end{proof}
\begin{definition}[acceleration In Meters Per Second Squared]
  \label{def:physics:phyx-mini-0697:phyxminiproblems-problemphyxmini0697-accelerationinmeterspersecondsquared}
  \lean{PhyXMiniProblems.ProblemPhyXMini0697.accelerationInMetersPerSecondSquared}
  \uses{def:physics:phyx-mini-0697:phyxminiproblems-problemphyxmini0697-accelerationmagnitudequantity}
  Metre-per-second-squared readout of tangential-acceleration magnitude
\end{definition}
\begin{proof}
  This is the declared model definition of acceleration In Meters Per Second Squared.
\end{proof}
\begin{definition}[Thin Uniform Rod]
  \label{def:physics:phyx-mini-0697:phyxminiproblems-problemphyxmini0697-thinuniformrod}
  \lean{PhyXMiniProblems.ProblemPhyXMini0697.ThinUniformRod}
  \uses{def:physics:phyx-mini-0697:phyxminiproblems-problemphyxmini0697-massquantity, def:physics:phyx-mini-0697:phyxminiproblems-problemphyxmini0697-lengthquantity, def:physics:phyx-mini-0697:phyxminiproblems-problemphyxmini0697-linearmassdensityquantity}
  A thin uniform rod together with its independent center-of-mass observable
\end{definition}
\begin{proof}
  This is the declared model definition of Thin Uniform Rod.
\end{proof}
\begin{definition}[Figure Feature]
  \label{def:physics:phyx-mini-0697:phyxminiproblems-problemphyxmini0697-figurefeature}
  \lean{PhyXMiniProblems.ProblemPhyXMini0697.FigureFeature}
  Features explicitly visible or printed in the supplied image
\end{definition}
\begin{proof}
  This is the declared model definition of Figure Feature.
\end{proof}
\begin{definition}[Rod Figure Orientation]
  \label{def:physics:phyx-mini-0697:phyxminiproblems-problemphyxmini0697-rodfigureorientation}
  \lean{PhyXMiniProblems.ProblemPhyXMini0697.RodFigureOrientation}
  Orientation of the rod relative to the Cartesian axes in the image
\end{definition}
\begin{proof}
  This is the declared model definition of Rod Figure Orientation.
\end{proof}
\begin{definition}[Uniform Rod Figure]
  \label{def:physics:phyx-mini-0697:phyxminiproblems-problemphyxmini0697-uniformrodfigure}
  \lean{PhyXMiniProblems.ProblemPhyXMini0697.UniformRodFigure}
  \uses{def:physics:phyx-mini-0697:phyxminiproblems-problemphyxmini0697-figurefeature, def:physics:phyx-mini-0697:phyxminiproblems-problemphyxmini0697-rodfigureorientation}
  Literal geometric and scalar data transcribed from the primary image. Cell values are SI coordinate readouts of the schematic x, dx, and dm; the underlying rod quantities remain dimensionful
\end{definition}
\begin{proof}
  This is the declared model definition of Uniform Rod Figure.
\end{proof}
\begin{definition}[Rotation Axis]
  \label{def:physics:phyx-mini-0697:phyxminiproblems-problemphyxmini0697-rotationaxis}
  \lean{PhyXMiniProblems.ProblemPhyXMini0697.RotationAxis}
  Axis about which the numerical rod rotates
\end{definition}
\begin{proof}
  This is the declared model definition of Rotation Axis.
\end{proof}
\begin{definition}[Rotating Rod Experiment]
  \label{def:physics:phyx-mini-0697:phyxminiproblems-problemphyxmini0697-rotatingrodexperiment}
  \lean{PhyXMiniProblems.ProblemPhyXMini0697.RotatingRodExperiment}
  \uses{def:physics:phyx-mini-0697:phyxminiproblems-problemphyxmini0697-angularaccelerationmagnitudequantity, def:physics:phyx-mini-0697:phyxminiproblems-problemphyxmini0697-accelerationmagnitudequantity, def:physics:phyx-mini-0697:phyxminiproblems-problemphyxmini0697-thinuniformrod, def:physics:phyx-mini-0697:phyxminiproblems-problemphyxmini0697-uniformrodfigure, def:physics:phyx-mini-0697:phyxminiproblems-problemphyxmini0697-rotationaxis}
  The physical rotating-rod experiment. Tangential acceleration is a dimensionful observable at every SI coordinate, rather than a single scalar attached to the entire rod. Its governing law is imposed only on coordinates inside the material interval
\end{definition}
\begin{proof}
  This is the declared model definition of Rotating Rod Experiment.
\end{proof}
\begin{definition}[Matches Primary Rod Figure]
  \label{def:physics:phyx-mini-0697:phyxminiproblems-problemphyxmini0697-matchesprimaryrodfigure}
  \lean{PhyXMiniProblems.ProblemPhyXMini0697.MatchesPrimaryRodFigure}
  \uses{def:physics:phyx-mini-0697:phyxminiproblems-problemphyxmini0697-massinkilograms, def:physics:phyx-mini-0697:phyxminiproblems-problemphyxmini0697-lengthinmeters, def:physics:phyx-mini-0697:phyxminiproblems-problemphyxmini0697-rotatingrodexperiment}
  Primary-image transcription. The last field is exactly the small-cell law dm = (M/L) dx; it records source evidence about uniformity and contains no center-of-mass or acceleration conclusion
\end{definition}
\begin{proof}
  This is the declared model definition of Matches Primary Rod Figure.
\end{proof}
\begin{definition}[Obeys Thin Uniform Rod Mass Model]
  \label{def:physics:phyx-mini-0697:phyxminiproblems-problemphyxmini0697-obeysthinuniformrodmassmodel}
  \lean{PhyXMiniProblems.ProblemPhyXMini0697.ObeysThinUniformRodMassModel}
  \uses{def:physics:phyx-mini-0697:phyxminiproblems-problemphyxmini0697-massinkilograms, def:physics:phyx-mini-0697:phyxminiproblems-problemphyxmini0697-lengthinmeters, def:physics:phyx-mini-0697:phyxminiproblems-problemphyxmini0697-linearmassdensityinkilogramspermeter, def:physics:phyx-mini-0697:phyxminiproblems-problemphyxmini0697-thinuniformrod}
  The continuum model for a uniform thin rod on [0,L]. It states positivity, constant density, recovery of total mass by integration, and the first-moment characterization of the independent center-of-mass observable. It does not state that the center is at L/2
\end{definition}
\begin{proof}
  This is the declared model definition of Obeys Thin Uniform Rod Mass Model.
\end{proof}
\begin{definition}[Matches Rotating Rod Problem]
  \label{def:physics:phyx-mini-0697:phyxminiproblems-problemphyxmini0697-matchesrotatingrodproblem}
  \lean{PhyXMiniProblems.ProblemPhyXMini0697.MatchesRotatingRodProblem}
  \uses{def:physics:phyx-mini-0697:phyxminiproblems-problemphyxmini0697-lengthinmeters, def:physics:phyx-mini-0697:phyxminiproblems-problemphyxmini0697-angularaccelerationinradianspersecondsquared, def:physics:phyx-mini-0697:phyxminiproblems-problemphyxmini0697-rotatingrodexperiment}
  Only the printed numerical data and stated rotation axis are recorded here. In particular, the source does not select a requested endpoint or assert a single tangential acceleration for the entire rod
\end{definition}
\begin{proof}
  This is the declared model definition of Matches Rotating Rod Problem.
\end{proof}
\begin{definition}[Obeys Rigid Rotation Kinematics]
  \label{def:physics:phyx-mini-0697:phyxminiproblems-problemphyxmini0697-obeysrigidrotationkinematics}
  \lean{PhyXMiniProblems.ProblemPhyXMini0697.ObeysRigidRotationKinematics}
  \uses{def:physics:phyx-mini-0697:phyxminiproblems-problemphyxmini0697-lengthinmeters, def:physics:phyx-mini-0697:phyxminiproblems-problemphyxmini0697-angularaccelerationinradianspersecondsquared, def:physics:phyx-mini-0697:phyxminiproblems-problemphyxmini0697-accelerationinmeterspersecondsquared, def:physics:phyx-mini-0697:phyxminiproblems-problemphyxmini0697-rotatingrodexperiment}
  Pointwise rigid-rotation kinematics. When the axis is through the center of mass, the radius of material coordinate x is |x - xCM|, and the standard magnitude law is aₜ = α r. No numerical endpoint value occurs in this law
\end{definition}
\begin{proof}
  This is the declared model definition of Obeys Rigid Rotation Kinematics.
\end{proof}
\begin{definition}[Has Position Independent Tangential Acceleration]
  \label{def:physics:phyx-mini-0697:phyxminiproblems-problemphyxmini0697-haspositionindependenttangentialacceleration}
  \lean{PhyXMiniProblems.ProblemPhyXMini0697.HasPositionIndependentTangentialAcceleration}
  \uses{def:physics:phyx-mini-0697:phyxminiproblems-problemphyxmini0697-lengthinmeters, def:physics:phyx-mini-0697:phyxminiproblems-problemphyxmini0697-accelerationinmeterspersecondsquared, def:physics:phyx-mini-0697:phyxminiproblems-problemphyxmini0697-rotatingrodexperiment}
  This predicate expresses the literal but physically inappropriate reading that one tangential-acceleration value applies at every point of the rod. Its negation is a derived conclusion for the nonzero-angular-acceleration experiment
\end{definition}
\begin{proof}
  This is the declared model definition of Has Position Independent Tangential Acceleration.
\end{proof}
\begin{lemma}[thin Uniform Rod center Of Mass]
  \label{lem:physics:phyx-mini-0697:phyxminiproblems-problemphyxmini0697-thinuniformrod-centerofmass}
  \lean{PhyXMiniProblems.ProblemPhyXMini0697.thinUniformRod_centerOfMass}
  \uses{def:physics:phyx-mini-0697:phyxminiproblems-problemphyxmini0697-lengthinmeters, def:physics:phyx-mini-0697:phyxminiproblems-problemphyxmini0697-thinuniformrod, def:physics:phyx-mini-0697:phyxminiproblems-problemphyxmini0697-obeysthinuniformrodmassmodel}
  Uniform density and the first-moment law put the center at L/2
\end{lemma}
\begin{proof}
  The conclusion follows from the governing relations and figure readouts named in the statement of thin Uniform Rod center Of Mass.
\end{proof}
\begin{lemma}[numerical Rod center Of Mass]
  \label{lem:physics:phyx-mini-0697:phyxminiproblems-problemphyxmini0697-numericalrod-centerofmass}
  \lean{PhyXMiniProblems.ProblemPhyXMini0697.numericalRod_centerOfMass}
  \uses{def:physics:phyx-mini-0697:phyxminiproblems-problemphyxmini0697-lengthinmeters, def:physics:phyx-mini-0697:phyxminiproblems-problemphyxmini0697-rotatingrodexperiment, def:physics:phyx-mini-0697:phyxminiproblems-problemphyxmini0697-obeysthinuniformrodmassmodel, def:physics:phyx-mini-0697:phyxminiproblems-problemphyxmini0697-matchesrotatingrodproblem}
  For the stated 1.60 m rod, the midpoint is 0.80 m
\end{lemma}
\begin{proof}
  The conclusion follows from the governing relations and figure readouts named in the statement of numerical Rod center Of Mass.
\end{proof}
\begin{lemma}[numerical Rod tangential Acceleration Profile]
  \label{lem:physics:phyx-mini-0697:phyxminiproblems-problemphyxmini0697-numericalrod-tangentialaccelerationprofile}
  \lean{PhyXMiniProblems.ProblemPhyXMini0697.numericalRod_tangentialAccelerationProfile}
  \uses{def:physics:phyx-mini-0697:phyxminiproblems-problemphyxmini0697-lengthinmeters, def:physics:phyx-mini-0697:phyxminiproblems-problemphyxmini0697-accelerationinmeterspersecondsquared, def:physics:phyx-mini-0697:phyxminiproblems-problemphyxmini0697-rotatingrodexperiment, def:physics:phyx-mini-0697:phyxminiproblems-problemphyxmini0697-obeysthinuniformrodmassmodel, def:physics:phyx-mini-0697:phyxminiproblems-problemphyxmini0697-matchesrotatingrodproblem, def:physics:phyx-mini-0697:phyxminiproblems-problemphyxmini0697-obeysrigidrotationkinematics}
  The acceleration profile is 6 * |x - 0.8| m/s² on the rod
\end{lemma}
\begin{proof}
  The conclusion follows from the governing relations and figure readouts named in the statement of numerical Rod tangential Acceleration Profile.
\end{proof}
\begin{lemma}[endpoint tangential Accelerations]
  \label{lem:physics:phyx-mini-0697:phyxminiproblems-problemphyxmini0697-endpoint-tangentialaccelerations}
  \lean{PhyXMiniProblems.ProblemPhyXMini0697.endpoint_tangentialAccelerations}
  \uses{def:physics:phyx-mini-0697:phyxminiproblems-problemphyxmini0697-lengthinmeters, def:physics:phyx-mini-0697:phyxminiproblems-problemphyxmini0697-accelerationinmeterspersecondsquared, def:physics:phyx-mini-0697:phyxminiproblems-problemphyxmini0697-rotatingrodexperiment, def:physics:phyx-mini-0697:phyxminiproblems-problemphyxmini0697-obeysthinuniformrodmassmodel, def:physics:phyx-mini-0697:phyxminiproblems-problemphyxmini0697-matchesrotatingrodproblem, def:physics:phyx-mini-0697:phyxminiproblems-problemphyxmini0697-obeysrigidrotationkinematics}
  Both endpoints have tangential-acceleration magnitude 4.8 m/s²
\end{lemma}
\begin{proof}
  The conclusion follows from the governing relations and figure readouts named in the statement of endpoint tangential Accelerations.
\end{proof}
\begin{lemma}[center tangential Acceleration]
  \label{lem:physics:phyx-mini-0697:phyxminiproblems-problemphyxmini0697-center-tangentialacceleration}
  \lean{PhyXMiniProblems.ProblemPhyXMini0697.center_tangentialAcceleration}
  \uses{def:physics:phyx-mini-0697:phyxminiproblems-problemphyxmini0697-lengthinmeters, def:physics:phyx-mini-0697:phyxminiproblems-problemphyxmini0697-accelerationinmeterspersecondsquared, def:physics:phyx-mini-0697:phyxminiproblems-problemphyxmini0697-rotatingrodexperiment, def:physics:phyx-mini-0697:phyxminiproblems-problemphyxmini0697-obeysthinuniformrodmassmodel, def:physics:phyx-mini-0697:phyxminiproblems-problemphyxmini0697-matchesrotatingrodproblem, def:physics:phyx-mini-0697:phyxminiproblems-problemphyxmini0697-obeysrigidrotationkinematics}
  The point on the rotation axis has zero tangential acceleration
\end{lemma}
\begin{proof}
  The conclusion follows from the governing relations and figure readouts named in the statement of center tangential Acceleration.
\end{proof}
\begin{lemma}[tangential Acceleration le endpoint Value]
  \label{lem:physics:phyx-mini-0697:phyxminiproblems-problemphyxmini0697-tangentialacceleration-le-endpointvalue}
  \lean{PhyXMiniProblems.ProblemPhyXMini0697.tangentialAcceleration_le_endpointValue}
  \uses{def:physics:phyx-mini-0697:phyxminiproblems-problemphyxmini0697-lengthinmeters, def:physics:phyx-mini-0697:phyxminiproblems-problemphyxmini0697-accelerationinmeterspersecondsquared, def:physics:phyx-mini-0697:phyxminiproblems-problemphyxmini0697-rotatingrodexperiment, def:physics:phyx-mini-0697:phyxminiproblems-problemphyxmini0697-obeysthinuniformrodmassmodel, def:physics:phyx-mini-0697:phyxminiproblems-problemphyxmini0697-matchesrotatingrodproblem, def:physics:phyx-mini-0697:phyxminiproblems-problemphyxmini0697-obeysrigidrotationkinematics}
  4.8 m/s² is an upper bound attained at both rod endpoints
\end{lemma}
\begin{proof}
  The conclusion follows from the governing relations and figure readouts named in the statement of tangential Acceleration le endpoint Value.
\end{proof}
\begin{lemma}[tangential Acceleration not position Independent]
  \label{lem:physics:phyx-mini-0697:phyxminiproblems-problemphyxmini0697-tangentialacceleration-not-positionindependent}
  \lean{PhyXMiniProblems.ProblemPhyXMini0697.tangentialAcceleration_not_positionIndependent}
  \uses{def:physics:phyx-mini-0697:phyxminiproblems-problemphyxmini0697-rotatingrodexperiment, def:physics:phyx-mini-0697:phyxminiproblems-problemphyxmini0697-obeysthinuniformrodmassmodel, def:physics:phyx-mini-0697:phyxminiproblems-problemphyxmini0697-matchesrotatingrodproblem, def:physics:phyx-mini-0697:phyxminiproblems-problemphyxmini0697-obeysrigidrotationkinematics, def:physics:phyx-mini-0697:phyxminiproblems-problemphyxmini0697-haspositionindependenttangentialacceleration}
  The rod does not have one position-independent tangential acceleration
\end{lemma}
\begin{proof}
  The conclusion follows from the governing relations and figure readouts named in the statement of tangential Acceleration not position Independent.
\end{proof}
\begin{definition}[Answer Choice]
  \label{def:physics:phyx-mini-0697:phyxminiproblems-problemphyxmini0697-answerchoice}
  \lean{PhyXMiniProblems.ProblemPhyXMini0697.AnswerChoice}
  Labels of the four answers printed with the problem
\end{definition}
\begin{proof}
  This is the declared model definition of Answer Choice.
\end{proof}
\begin{definition}[tangential Acceleration In Meters Per Second Squared]
  \label{def:physics:phyx-mini-0697:phyxminiproblems-problemphyxmini0697-answerchoice-tangentialaccelerationinmeterspersecondsquared}
  \lean{PhyXMiniProblems.ProblemPhyXMini0697.AnswerChoice.tangentialAccelerationInMetersPerSecondSquared}
  \uses{def:physics:phyx-mini-0697:phyxminiproblems-problemphyxmini0697-answerchoice}
  SI tangential-acceleration number printed beside each answer label
\end{definition}
\begin{proof}
  This is the declared model definition of tangential Acceleration In Meters Per Second Squared.
\end{proof}
\begin{definition}[recorded Answer Choice]
  \label{def:physics:phyx-mini-0697:phyxminiproblems-problemphyxmini0697-recordedanswerchoice}
  \lean{PhyXMiniProblems.ProblemPhyXMini0697.recordedAnswerChoice}
  \uses{def:physics:phyx-mini-0697:phyxminiproblems-problemphyxmini0697-answerchoice}
  Answer label recorded in the supplied dataset
\end{definition}
\begin{proof}
  This is the declared model definition of recorded Answer Choice.
\end{proof}
\begin{definition}[Matches Displayed Tangential Acceleration At Position]
  \label{def:physics:phyx-mini-0697:phyxminiproblems-problemphyxmini0697-matchesdisplayedtangentialaccelerationatposition}
  \lean{PhyXMiniProblems.ProblemPhyXMini0697.MatchesDisplayedTangentialAccelerationAtPosition}
  \uses{def:physics:phyx-mini-0697:phyxminiproblems-problemphyxmini0697-accelerationinmeterspersecondsquared, def:physics:phyx-mini-0697:phyxminiproblems-problemphyxmini0697-rotatingrodexperiment, def:physics:phyx-mini-0697:phyxminiproblems-problemphyxmini0697-answerchoice}
  A displayed number matches the physical acceleration at a specified material coordinate. Requiring the coordinate prevents answer metadata from silently selecting an endpoint
\end{definition}
\begin{proof}
  This is the declared model definition of Matches Displayed Tangential Acceleration At Position.
\end{proof}
% --- Archon declaration topology end ---
% --- Archon physics formalization source end ---
